\section{Sharp late coefficients of a resonant Lambert inverse}
\label{lm:sec:main}

The newest incoming Lambert manuscript leaves a precise question about the
large-order coefficients of its inverse, under the heading ``Sharp
large-order constants of the inverse sectors'' \cite{RepoLambert}.
We prove its proposed even and odd constants, obtain a further relative
correction, and identify the corresponding logarithm and dilogarithm terms
inside the initial Borel disk.  The proof also explains which changes to
the analytic core or to the finite forward prefix can affect those terms.
General algebras of factorially divergent series provide a broader setting
for such questions \cite{Borinsky}; the argument below supplies direct
estimates for the specific coefficient extraction and its full sum over
the number of factors.  Lagrange inversion is used in its usual formal
coefficient sense \cite{Gessel}.

\subsection{The source equation and its normalization}
\label{lm:subsec:normalization}

Set $A=4\pi^2$ and consider the double Lambert sum
\[
 \mathcal F(t)=\sum_{a,b\geq1}\frac{e^{-abt}}{ab}
              =\sum_{n\geq1}\frac{d(n)}{n}e^{-nt},\qquad t>0.
\]
Here $d(n)$ counts the positive divisors of $n$. We use the Stieltjes
convention $\zeta(1+w)=w^{-1}+\sum_{j\ge0}(-1)^j\gamma_jw^j/j!$.
Its Mellin integrand is $\Gamma(w)\zeta(1+w)^2t^{-w}$.  The triple pole
at zero gives the core
\[
 \frac12\bigl(\log(1/t)+\gamma\bigr)^2+C_* ,
 \qquad C_* =\frac{\zeta(2)}2-\gamma^2-2\gamma_1.
\]
The negative integer residues give the formal tail
\begin{equation}\label{lm:eq:forward}
 R(q)=\sum_{n\geq1}r_nq^n,\qquad
 r_1=-\frac14,\qquad
 r_{2m}=\frac{4\Gamma(2m)}{(2m)A^{2m}}\zeta(2m)^2,
 \qquad r_{2m+1}=0\quad(m\geq1).
\end{equation}
In particular $r_2=1/288$.  The real completion of this expansion is
treated in the predecessor \cite{RepoLambert}.  Here we need only the
explicit formal coefficients in \eqref{lm:eq:forward}.

For $Y$ large, the exact-core variables of that manuscript are
$X=\sqrt{2(Y-C_*)}$, $u=X^{-1}$, $q=e^{\gamma-X}$, and $t=qe^{-V}$.
Their formal inverse equation is
\begin{equation}\label{lm:eq:inverse}
 V+\frac u2V^2+uR(qe^{-V})=0,
 \qquad V(q,u)=\sum_{n\geq1}C_n(u)q^n.
\end{equation}
Throughout the coefficient analysis, $q$ and $u$ are independent; the
physical specialization is a later operation.  The formal implicit
function theorem gives a unique solution.  The predecessor proves that
each $C_n$ is a polynomial and obtains
\begin{equation}\label{lm:eq:finite}
 C_n(u)=-\sum_{k=1}^n\frac{u^k}{k!}[q^n]R(q)^k
 \sum_{j=0}^{k-1}\binom{k-1}{j}n^{k-1-j}(k)_j\left(\frac u2\right)^j.
\end{equation}
Its conjecture is that the dominant configurations contain one large
even index and $k-1$ linear indices.  The issue left open there is the
uniform justification of summing that approximation over $k$.

\begin{theorem}[Sharp Lambert inverse amplitudes]\label{lm:thm:sharp}
For every $U>0$, uniformly for $|u|\leq U$,
\begin{equation}\label{lm:eq:sharp}
 C_n(u)=-\frac{2u\Gamma(n)}{nA^n}
 \left(e^{-Au/4}+(-1)^ne^{Au/4}\right)
 +O_U\!\left(\frac{\Gamma(n)}{n^2A^n}\right).
\end{equation}
Consequently, for every fixed real $u>0$,
\begin{align}
 C_{2m}(u)&\sim-\frac{4u\Gamma(2m)}{(2m)A^{2m}}
                         \cosh(Au/4),\label{lm:eq:even}\\
 C_{2m+1}(u)&\sim\frac{4u\Gamma(2m+1)}{(2m+1)A^{2m+1}}
                         \sinh(Au/4).\label{lm:eq:odd}
\end{align}
These are the two conjectured constants of the incoming manuscript.
\end{theorem}

The additive formulation \eqref{lm:eq:sharp} also covers complex values
at which one of the displayed hyperbolic functions vanishes.  It is
stronger than a separate positive-real ratio asymptotic.

\subsection{Analytic-core universality}
\label{lm:subsec:universality}

\begin{theorem}[Universality of the leading action amplitudes]
\label{lm:thm:universal}
Let $A>0$, $r_1=c$, and suppose
\begin{equation}\label{lm:eq:tail-assumption}
 r_n=\frac{\Gamma(n)}{nA^n}
       \bigl(\sigma_++(-1)^n\sigma_-\bigr)
       +O\!\left(\frac{\Gamma(n)}{n^2A^n}\right),\qquad n\geq2.
\end{equation}
Let $\Psi(v,u)$ be analytic near $v=0$ and $u$ in an open parameter
domain, with $\Psi(0,u)=0$ and $\partial_v\Psi(0,u)=1$; assume a
common neighborhood of $v=0$ on each compact parameter set.  The unique
formal solution
\[
 \Psi(V,u)+uR(qe^{-V})=0,
 \qquad V(q,u)=\sum_{n\geq1}C_n(u)q^n,
\]
satisfies, locally uniformly in $u$,
\begin{equation}\label{lm:eq:universal}
 C_n(u)=-\frac{u\Gamma(n)}{nA^n}
 \left(\sigma_+e^{Acu}+(-1)^n\sigma_-e^{-Acu}\right)
 +O\!\left(\frac{\Gamma(n)}{n^2A^n}\right).
\end{equation}
Thus nonlinear changes to the analytic core and finite changes to
$r_2,r_3,\ldots$ preserve the leading amplitudes.  Changing $r_1=c$
multiplies the two amplitudes by the corresponding exponential factors.
\end{theorem}

\begin{lemma}[Factorial convolutions]\label{lm:lem:factorials}
For integers $N\geq j\geq1$,
\begin{equation}\label{lm:eq:factorial-convolution}
 \sum_{\substack{\ell_1+\cdots+\ell_j=N\\\ell_i\geq1}}
       \prod_{i=1}^j(\ell_i-1)!
 \leq 3^{j-1}(N-j)!.
\end{equation}
Suppose $A|r_1|\leq K$ and $A^m|r_m|\leq K(m-2)!$ for $m\geq2$.
For $1\leq k<n$, with $N=n-k$,
\begin{equation}\label{lm:eq:all-configurations}
 A^n|[q^n]R(q)^k|\leq(4K)^k\Gamma(N).
\end{equation}
Moreover
\begin{equation}\label{lm:eq:one-large}
 [q^n]R(q)^k=kc^{k-1}r_{n-k+1}+E_{n,k},
 \qquad A^n|E_{n,k}|\leq(4K)^k\Gamma(N-1)\quad(N\geq2).
\end{equation}
\end{lemma}

\begin{proof}
The elementary inequality
\[
 \sum_{p=0}^M p!(M-p)!\leq3M!
\]
holds because the two endpoint reciprocal binomial coefficients sum to
$2$, while each interior one is at most $1/M$; the cases $M=0,1$ are
immediate.  Induction gives \eqref{lm:eq:factorial-convolution}.

In a composition $m_1+\cdots+m_k=n$ put $\ell_i=m_i-1$ and let $j$
be the number of positive $\ell_i$.  The absolute sum is at most
\[
 K^k\sum_{j=1}^{\min(k,N)}\binom{k}{j}3^{j-1}(N-j)!.
\]
Bounding each factorial by $(N-1)!$ proves
\eqref{lm:eq:all-configurations}.  The terms with $j=1$ give exactly
$kc^{k-1}r_{n-k+1}$; for the remaining $j\geq2$ one may instead bound
by $(N-2)!$, proving \eqref{lm:eq:one-large}.
\end{proof}

\begin{proof}[Proof of Theorem~\ref{lm:thm:universal}]
Write $G(v,u)=v/\Psi(v,u)$, analytically continued at $v=0$, so that
$G(0,u)=1$.  Lagrange inversion in an auxiliary multiplier gives
\begin{equation}\label{lm:eq:general-finite}
 C_n(u)=-\sum_{k=1}^n\frac{u^k}{k!}n^{k-1}[q^n]R^k H_{n,k}(u),
\end{equation}
where
\begin{align}
 H_{n,k}(u)&=(-n)^{1-k}
  \left.\partial_v^{k-1}\{e^{-nv}G(v,u)^k\}\right|_{v=0}\notag\\
 &=\sum_{j=0}^{k-1}\binom{k-1}{j}(-n)^{-j}
              \left.\partial_v^jG(v,u)^k\right|_{v=0}.
 \label{lm:eq:H}
\end{align}
Only $k\leq n$ can contribute, so the auxiliary multiplier may be set
to one formally.  This also proves the coefficient formula without
assuming convergence of $R$.

Fix a compact parameter set.  Choose $\rho>0$ and $M\geq1$ such that
$|G(v,u)|\leq M$ on $|v|\leq\rho$.  Cauchy's estimate yields
$|\partial_v^jG^k(0,u)|\leq j!M^k\rho^{-j}$.  Since
$j!\leq n^j$ for $j\leq n$, there is a constant $B$ with
$|H_{n,k}|\leq B^k$.  The sharper inequality
$j!\leq jn^{j-1}$ for $j\geq1$ gives
\begin{equation}\label{lm:eq:H-error}
 |H_{n,k}-1|\leq
 \frac{M^k(k-1)\rho^{-1}(1+\rho^{-1})^{k-2}}n
 \leq \frac{kB_1^k}{n}.
\end{equation}
Here and below constants may depend on the compact parameter set.

Assumption \eqref{lm:eq:tail-assumption} supplies the constant $K$ of
Lemma~\ref{lm:lem:factorials}.  After division by
$\Gamma(n)/(nA^n)$, the absolute $k$th summand of
\eqref{lm:eq:general-finite}, for $k<n$, is bounded by
\begin{equation}\label{lm:eq:k-majorant}
 \frac{D^k}{k!}\frac{n^k\Gamma(n-k)}{\Gamma(n)}.
\end{equation}
For $k\leq n/2$ the last ratio is at most $2^k$, since it is the
product of $n/(n-j)$ for $1\leq j\leq k$.  This produces the
summable majorant $(2D)^k/k!$.

For $n/2<k<n$, use $\Gamma(n-k)\leq\Gamma(\lceil n/2\rceil)$ and
$n^k/k!\leq(en/k)^k\leq(2e)^k$.  The entire contribution of this
range is bounded by
\begin{equation}\label{lm:eq:large-k}
 n\max(1,2eD)^n\frac{\Gamma(\lceil n/2\rceil)}{\Gamma(n)},
\end{equation}
which is smaller than every inverse power of $n$.  For $k=n$ one
uses $[q^n]R^n=c^n$ directly; the normalized bound is a constant to
the $n$th power divided by $\Gamma(n)$ and has the same property.

It remains to sharpen the range $k\leq n/2$.  By
\eqref{lm:eq:one-large}, the normalized total contribution of all
$E_{n,k}$ is $O(1/n)$: indeed its $k$th bound differs from
\eqref{lm:eq:k-majorant} by the extra factor
$(n-k-1)^{-1}=O(1/n)$.  For the remaining term put $m=n-k+1$.
The gamma ratio needed in its evaluation is
\begin{equation}\label{lm:eq:gamma-ratio}
 B_{n,k}:=\frac{n^k\Gamma(n-k+1)}{(n-k+1)\Gamma(n)}
 =\frac n{n-k+1}\prod_{j=1}^{k-1}\frac n{n-j}.
\end{equation}
For $k\leq n/2$ it is at most $2^k$, and
$|B_{n,k}-1|\leq Ck^2 2^k/n$.  To see the latter, take the logarithm
of the finite product, use $-\log(1-x)\leq2x$ for $0\leq x\leq1/2$,
and then $e^y-1\leq ye^y$.

The error in \eqref{lm:eq:tail-assumption}, with $m\geq n/2$, is
also $O(1/n)$ after this normalization.  Combining it with
\eqref{lm:eq:H-error} and \eqref{lm:eq:gamma-ratio} shows that the
$k$th term may be replaced, with a total summable error $O(1/n)$, by
\[
 -\frac{u(Acu)^{k-1}}{(k-1)!}
       \bigl(\sigma_++(-1)^{n-k+1}\sigma_-\bigr).
\]
All error bounds have a polynomial in $k$ times a constant to the
$k$th power divided by $(k-1)!$.  Summing the two exponential series
proves \eqref{lm:eq:universal}, uniformly on the compact parameter set.
\end{proof}

\begin{proof}[Proof of Theorem~\ref{lm:thm:sharp}]
In \eqref{lm:eq:forward}, $\zeta(n)^2=1+O(2^{-n})$ for $n\geq2$.
Thus \eqref{lm:eq:tail-assumption} holds with $c=-1/4$ and
$\sigma_+=\sigma_-=2$.  Apply Theorem~\ref{lm:thm:universal} to
$\Psi(v,u)=v+uv^2/2$.  The identities for $\cosh$ and $\sinh$ give
\eqref{lm:eq:even}--\eqref{lm:eq:odd}.
\end{proof}

\subsection{The first correction and the analytic data it detects}
\label{lm:subsec:first}

\begin{theorem}[First relative correction]\label{lm:thm:first}
For the Lambert inverse \eqref{lm:eq:inverse}, put
\begin{equation}\label{lm:eq:theta-H}
 \theta=Au/4,\qquad
 P(u)=\frac{1+u}{2}\theta^2+\frac{uA^2}{288},\qquad
 D_n(u)=\frac{nA^nC_n(u)}{4\Gamma(n)}.
\end{equation}
Let $E_n(\theta)=\cosh\theta$ for even $n$ and
$E_n(\theta)=-\sinh\theta$ for odd $n$.  Uniformly for $|u|\leq U$,
\begin{equation}\label{lm:eq:first}
 D_n(u)=-uE_n(\theta)
 -\frac un\left(P(u)E_n(\theta)+(u+2)\theta E_n'(\theta)\right)
 +O_U(n^{-2}).
\end{equation}
The correction detects the quadratic core and $r_2=1/288$, although
neither affects the leading amplitudes.
\end{theorem}

\begin{proof}
We retain one more order in the preceding proof.  Uniformly for
$n\geq8$ and $k\leq n/2$, with $\ell=k-1$, expansion of the logarithm of
\eqref{lm:eq:gamma-ratio} gives
\begin{equation}\label{lm:eq:gamma-first}
 B_{n,k}=1+\frac{\ell(\ell+3)}{2n}
             +O\!\left(\frac{k^4B_2^k}{n^2}\right).
\end{equation}
For example, the sum of the first-order logarithmic terms is
$\ell(\ell+3)/(2n)$; the sum of their squared arguments is
$O(k^3/n^2)$, and the exponential remainder is bounded by
$O(k^4B_2^k/n^2)$.

For the quadratic core, \eqref{lm:eq:H} similarly gives
\begin{equation}\label{lm:eq:H-first}
 H_{n,k}=1+\frac{uk(k-1)}{2n}
              +O_U\!\left(\frac{k^4B_3^k}{n^2}\right).
\end{equation}
This is a uniform estimate, not just a fixed-$k$ expansion.  In the
Cauchy argument, the terms $j\geq2$ are bounded using
$j!\leq j(j-1)n^{j-2}$; their binomial sum is bounded by a polynomial
in $k$ times a constant to the $k$th power, divided by $n^2$.
For a general core $\Psi(v,u)=v+b(u)v^2+O(v^3)$, the first correction
in \eqref{lm:eq:H-first} is $b(u)k(k-1)/n$.

The one-large-index configurations therefore contribute the operator
\[
 \frac12\left((1+u)(\theta\partial_\theta)^2
                    +(3+u)\theta\partial_\theta\right)
\]
applied to the parity exponential sum.  Since $E_n''=E_n$, its value is
\begin{equation}\label{lm:eq:operator-first}
 \frac{1+u}{2}\theta^2E_n+(u+2)\theta E_n'.
\end{equation}

There is one further contribution of normalized order $1/n$: exactly
two indices are at least two, one equals $2$, and the other equals
$n-k$.  For $k\leq n/2$ and sufficiently large $n$, its coefficient
is exactly
\[
 k(k-1)c^{k-2}r_2r_{n-k}.
\]
Using its leading gamma ratio and summing over $k$ gives
$-A^2r_2u^2E_n/n$.  Together with \eqref{lm:eq:operator-first}, this
is precisely the correction in \eqref{lm:eq:first}.

For completeness, the configurations omitted in this refinement have
a uniform $O(n^{-2})$ bound after normalization.  Those with $j\geq3$
indices at least two satisfy the strengthened bound
$(4K)^k\Gamma(n-k-2)$ in Lemma~\ref{lm:lem:factorials}.  With
$k\leq n/2$ the two additional inverse factors are $O(n^{-2})$.
For $j=2$, remove the two endpoint configurations just used.  If
$N=n-k$, both remaining increments are at least two, and
\begin{equation}\label{lm:eq:two-interior}
 \sum_{\ell=2}^{N-2}(\ell-1)!(N-\ell-1)!
       \leq4(N-3)!.
\end{equation}
Indeed, put $M=N-2$, divide by $M!$, and use
$\sum_{i=1}^{M-1}\binom Mi^{-1}\leq4/M$: the two endpoints
contribute $2/M$, while the remaining binomial coefficients are at
least $M(M-1)/2$.  Again two inverse powers of $n$ result.  The
gamma and derivative remainders in the retained configurations have
the same bound after summing their factorial majorants.  The range
$k>n/2$ is smaller than every inverse power by
\eqref{lm:eq:large-k}.  This proves the stated uniform remainder.
\end{proof}

\subsection{Exact logarithm--dilogarithm decomposition in the Borel disk}
\label{lm:subsec:borel}

The coefficient asymptotics have a direct polylogarithmic interpretation.
It only requires summing absolutely convergent series in the initial
disk and consequently does not require a continuation theorem for the
inverse Borel transform.

\begin{corollary}[The two nearest Borel singularities]
\label{lm:cor:logs}
For
\[
 \widehat V_u(\xi)=\sum_{n\geq1}\frac{C_n(u)}{\Gamma(n)}\xi^{n-1}
\]
there is, on $|\xi|<A$, the exact decomposition
\begin{equation}\label{lm:eq:log-decomposition}
 \widehat V_u(\xi)=\frac{2u}{\xi}
       \left(e^{-\theta}\log(1-\xi/A)
                     +e^{\theta}\log(1+\xi/A)\right)+J_u(\xi),
\end{equation}
where $J_u$ is holomorphic in the disk and continuous on its closure,
locally uniformly in $u$.  The factor $1/\xi$ has a removable
singularity at zero in the displayed combination.  For every complex
$u\ne0$, both endpoints $+A$ and $-A$ are logarithmically singular,
with radial logarithmic coefficients respectively
\begin{equation}\label{lm:eq:radial-amplitudes}
 \frac{2u}{A}e^{-\theta},\qquad -\frac{2u}{A}e^{\theta}.
\end{equation}
The corresponding radially regularized limits exist.  The Borel radius
is exactly $A$, and the ordinary $q$-series has radius zero.  At $u=0$
the formal solution is identically zero.
\end{corollary}

\begin{proof}
Divide \eqref{lm:eq:sharp} by $\Gamma(n)$ and sum using
$\sum_{n\geq1}z^n/n=-\log(1-z)$.  The error series is absolutely
and uniformly convergent on $|\xi|\leq A$, because its coefficients
are bounded by a constant times $A^{-n}/n^2$.  This proves
\eqref{lm:eq:log-decomposition}.  The coefficients in
\eqref{lm:eq:radial-amplitudes} never vanish for $u\ne0$, proving
both endpoint singularities and the exact Borel radius.  At least one
of $\cosh\theta$ and $\sinh\theta$ is nonzero, so
\eqref{lm:eq:sharp} on that parity subsequence also proves that the
ordinary radius is zero.  For $u=0$, formal uniqueness gives $V=0$.
\end{proof}

\begin{corollary}[An explicit dilogarithmic correction]
\label{lm:cor:dilog}
With $P$ from \eqref{lm:eq:theta-H}, set
\begin{equation}\label{lm:eq:Q}
 Q_+(u)=P(u)-(u+2)\theta,\qquad
 Q_-(u)=P(u)+(u+2)\theta.
\end{equation}
The following identity holds throughout $|\xi|<A$:
\begin{align}
 \widehat V_u(\xi)={}&\frac{2u}{\xi}
       \left(e^{-\theta}\log(1-\xi/A)
                         +e^{\theta}\log(1+\xi/A)\right)\notag\\
 &-\frac{2u}{\xi}
       \left(e^{-\theta}Q_+(u)\Li_2(\xi/A)
                         +e^{\theta}Q_-(u)\Li_2(-\xi/A)\right)
       +K_u(\xi).\label{lm:eq:dilog-decomposition}
\end{align}
Both $K_u$ and $\partial_\xi K_u$ extend continuously to the closed
disk, locally uniformly in $u$.
\end{corollary}

\begin{proof}
Theorem~\ref{lm:thm:first} is equivalent to
\begin{align*}
 \frac{C_n(u)}{\Gamma(n)}
 ={}&-\frac{2u}{nA^n}\bigl(e^{-\theta}+(-1)^ne^{\theta}\bigr)\\
 &-\frac{2u}{n^2A^n}
       \bigl(e^{-\theta}Q_+(u)+(-1)^ne^{\theta}Q_-(u)\bigr)
       +O_U(A^{-n}n^{-3}).
\end{align*}
Sum this identity with $\Li_2(z)=\sum_{n\geq1}z^n/n^2$.
The remainder and its once-differentiated series converge absolutely
and uniformly on $|\xi|\leq A$, proving the assertion.
\end{proof}

These are statements from inside the initial Borel disk.  They neither
assert analytic continuation through that circle nor identify all
subsequent singularities.  They do distinguish the two leading endpoint
amplitudes, improve the predecessor's radius-only result to every
nonzero complex $u$, and express the first correction by explicit
classical polylogarithms.
